\section{Reproduction Entry Points}
\label{sec:reproduction}

Run commands from \texttt{t-i/}. The reportable Runpod workflow is:
\par\noindent
\begin{minipage}{\linewidth}
\scriptsize\ttfamily
bash scripts/runpod\_pusht.sh setup\\
bash scripts/runpod\_pusht.sh prepare\\
bash scripts/runpod\_pusht.sh smoke\\
bash scripts/runpod\_pusht.sh train\\
RIPL\_RUN\_DIR=/workspace/ripl-artifacts/runs/RUN \\
\hspace*{1em}bash scripts/runpod\_pusht.sh eval
\end{minipage}

The canonical configuration is
\path{configs/pusht_rgb_delta_pose.yaml}. The preparation stage invokes
\path{scripts/prepare_pusht_delta_pose_demos.sh}, applies the audited alignment through
\path{scripts/replay_trajectory_aligned.py}, and validates the result with
\path{scripts/validate_pusht_replay.py}.

Resume only from a validated periodic checkpoint:
\par\noindent
\begin{minipage}{\linewidth}
\scriptsize\ttfamily
RIPL\_RESUME=/path/to/step\_15000.pt \\
\hspace*{1em}bash scripts/runpod\_pusht.sh train \\
\hspace*{2em}--total-iters 50000
\end{minipage}

The final evaluator loads \path{checkpoints/best_success_once.pt} for seeds 0, 1, and 2, saves
100 episodes per seed, and writes \path{evaluation-final/summary.json}. Exact setup details,
metric semantics, artifact requirements, and historical experiment provenance are documented in
\path{t-i/README.md} and \path{results/pusht_progress_2026-08-26.md}.

T-II tooling is independently testable from the repository root:
\par\noindent
\begin{minipage}{\linewidth}
\scriptsize\ttfamily
python -m pip install -e './t-i[dev]' -e './t-ii[dev]'\\
python -m pytest -q t-ii/tests
\end{minipage}

After provisioning a Linux NVIDIA GPU with working Vulkan graphics, the documented
\path{t-ii/scripts/run_pusht_t2.sh} stages are \texttt{discover}, \texttt{analyze}, and
\texttt{targeted}. The targeted stage must not be run until discovery evidence defines two
contiguous pose regimes. T-III's scaffold is similarly tested with
\texttt{python -m pytest -q t-iii/tests}; these tests do not constitute T-III experiments.
